% Generated draft: requires mathematical and source review.
\documentclass[11pt]{article}
\usepackage[margin=25mm]{geometry}
\usepackage{fontspec}
\setmainfont{Latin Modern Roman}
\usepackage{amsmath,amssymb,mathtools,unicode-math}
\setmathfont{Latin Modern Math}
\usepackage{enumitem,needspace,xurl,hyperref}
\setlength{\parindent}{0pt}
\setlength{\parskip}{5pt}
\setlength{\emergencystretch}{3em}
\newcommand{\N}{\mathbb{N}}
\newcommand{\Z}{\mathbb{Z}}
\newcommand{\Q}{\mathbb{Q}}
\newcommand{\R}{\mathbb{R}}
\newcommand{\C}{\mathbb{C}}
\newcommand{\F}{\mathbb{F}}
\newcommand{\A}{\mathbb{A}}
\newcommand{\PP}{\mathbb{P}}
\newcommand{\E}{\mathbb{E}}
\newcommand{\eps}{\varepsilon}
\newcommand{\dd}{\,\mathrm{d}}
\begin{document}
\section*{UnsolvedMath ID: 30006277}
\textbf{Source problem code:} catalogue entry
\subsection{Definitions and mathematical statement}
Let \(\pi\) denote the family called the universal elliptic curve in the supplied record. Its base, compactification, and precise moduli interpretation are not supplied. The requested invariants are denoted by
\[
\bigl\langle\Lambda\tau_k(0)\bigr\rangle^{\pi,d}_{g,1}.
\]
Here \(g\) is the genus, the subscript \(1\) indicates one marked point, and \(d\) is the degree parameter. The source does not specify the admissible ranges or stability conditions for these parameters, or the class \(\Lambda\).

The damaged displayed definition in the source can provisionally be reconstructed as
\[
\bigl\langle\Lambda\tau_k(0)\bigr\rangle^{\pi,d}_{g,1}
:=
\int_{[\mathcal M_{g,1}(\pi,d)]^{\mathrm{vir}}}
\operatorname{ft}^{*}\Lambda\,\psi_1^{k}\,\operatorname{ev}^{*}[0].
\]
In this reconstruction, \(\mathcal M_{g,1}(\pi,d)\) is a placeholder for the source's moduli space of maps associated with \(\pi\), and \([\mathcal M_{g,1}(\pi,d)]^{\mathrm{vir}}\) denotes its virtual fundamental class. The notation \(\psi_1\) conventionally denotes the first Chern class of the cotangent line at the marked point, \(\operatorname{ft}\) a forgetful or stabilization morphism, and \(\operatorname{ev}\) an evaluation morphism. The insertion \([0]\) appears to refer to the zero point or zero section of the elliptic target. These interpretations, including the pullback \(\operatorname{ft}^{*}\), require confirmation from the original correspondence omitted from the supplied excerpt. In particular, the moduli space, the morphisms and their targets, the meaning of \(\Lambda\), the meaning of \([0]\), and the dimension condition making the integral numerical have not been specified sufficiently to define the invariant uniquely.

\textbf{Problem.} With those data fixed as intended in the source, determine the invariants \(\langle\Lambda\tau_k(0)\rangle^{\pi,d}_{g,1}\) in general, including higher genus, for all admissible choices of the parameters and insertion. The source explicitly states that degeneration of the target elliptic curve to a rational curve is not allowed in this setting. The mathematical reason for this restriction is absent from the excerpt; it is retained as a qualification of the requested computation. This is a computational problem, and the supplied information does not specify a conjectural formula.
\subsection{Short English statement}
Compute the indicated descendant Gromov--Witten invariants of the universal elliptic curve in general, including higher genus, under the source's restriction that degeneration to a rational target is unavailable. The supplied excerpt omits definitions needed to specify these invariants uniquely, so the displayed integral is a provisional reconstruction requiring mathematical review.
\subsection{Sources}
\begin{itemize}
\item Supplied local UnsolvedMath registry record, permanent ID \(30006277\), source problem code \texttt{catalogue entry}, entitled \emph{Higher-Genus Gromov--Witten Invariants of Elliptic Targets}. The record belongs to the Oberwolfach Reports Open Problems collection and gives proposed year \(2025\); no individual proposer is supplied.
\item Supplied primary-source citation: \emph{Recent Trends in Algebraic Geometry} (2025), Oberwolfach Reports, pp.\ 1451--1516, DOI \texttt{10.4171/owr/2025/28}. \url{https://doi.org/10.4171/owr/2025/28}. The record describes Oberwolfach Reports as documenting workshops and their open-problem discussions; no precise page or contributor for this problem is identified.
\item Supplied literature reference: Iribar L\'opez, \emph{Noether--Lefschetz cycles}. \url{https://doi.org/10.1017/fmp.2025.10018}.
\item Supplied literature reference labelled ``original OWR report'': \url{https://doi.org/10.4171/owr/2025/28}.
\item Literature assessment dated exactly \texttt{2026-08-22}: the supplied curation labels the problem partially solved and states that the general all-genus computation remains unknown, while relative moduli constructions, degree calculations, and selected low-genus or special cases provide substantial but incomplete progress. This is supplied dated curation, not an independently verified assessment; its time-sensitive status and cited claims require independent verification.
\item The registry labels the statement \texttt{corrected\_verified} and reports that the official Oberwolfach report was checked, the question isolated from surrounding exposition, and its notation repaired for publication. These are supplied verification claims; no source visit or independent verification is asserted here.
\end{itemize}
\end{document}
